\section{Putting it together}\label{part:together}

This section assembles the moving parts into a single control loop: bounds and
their literals (\cref{part:foundations}), propagators that tighten them, branching
that commits values, and conflict analysis that learns from failure
(\cref{part:learning}). It then opens the loop's two busiest lines: what
``propagate to a fixed point'' actually runs, and in what order the SAT and CP
engines take turns (\cref{sec:propstep}); and how the next decision is
chosen once propagation stalls, where a deliberate design choice has CP-SAT drive
the search itself and lets a configurable decision strategy, by default led by SAT
branching, choose where to descend (\cref{sec:search}).

\subsection{The search loop, assembled}\label{sec:loop}
% Sources (VERIFIED):
%   ortools/sat/integer_search.cc -- main integer search loop: repeatedly
%     propagate, take a decision, on conflict analyze+backjump; return FEASIBLE
%     when the heuristic yields no decision (all variables fixed).
%   ortools/sat/sat_solver.cc -- SatSolver::Solve / Propagate / conflict handling.
%   ortools/sat/optimization.cc -- the objective-tightening wrapper around it.

Putting the pieces together, the engine is a single feasibility loop
(\srcla{integer\_search.cc}{1526}{SolveIntegerProblem}, entered through the
level-$0$ reset-and-assert wrapper
\srcla{integer\_search.cc}{1649}{ResetAndSolveIntegerProblem}) wrapped by a
thin optimizer that calls it repeatedly. The loop, \textsc{Solve}
(\cref{alg:loop}), is the heart:
every concept introduced so far appears in it exactly once, and each line points to
a section that expands it. It takes a set of \emph{assumptions} (literals forced
true for the duration of one call) and returns either a feasible solution or a
proof that none exists under them. The wrapper,
\textsc{MinimizeByBisection} (\cref{alg:bisection}), turns that feasibility engine into an optimizer: it
calls \textsc{Solve} under a trial bound on the objective and narrows that bound by
bisection until the optimum is pinned.

\begin{algorithm}[htbp]
\caption{The CP-SAT feasibility core, \textsc{Solve}}
\label{alg:loop}
\begin{algorithmic}[1]
\Procedure{Solve}{$A$} \Comment{$A$: assumption literals, forced true for this call}
  \LineComment{Prepare the model for a new run, with the assumptions enforced at the root level}
  \State reset to decision level $0$; assert every $\ell \in A$
  \Loop
    \LineComment{Run all propagators (cheapest first) until we reach a fixed point (\cref{sec:propstep})}
    \State \Call{PropagateToFixedPoint}{} \Comment{first SAT, then CP propagations}
    \If{conflict detected}
      \LineComment{Learn from the conflict and resolve it (\cref{sec:trace})}
      \If{current decision level $= 0$}
        \State \Return \textsc{Infeasible} \Comment{with the failing core $\subseteq A$, when $A \neq \emptyset$}
      \EndIf
      \State $c \gets \Call{AnalyzeConflict}{}$ \Comment{1-UIP nogood}
      \State \Call{Backjump}{\textsc{AssertionLevel}$(c)$}; \ \Call{Learn}{$c$}
    \Else
      \LineComment{Ask the decision strategies for a branching literal, or confirm none is left (\cref{sec:search})}
      \State $\ell \gets \Call{NextDecision}{}$ \Comment{default: (1)~VSIDS, (2)~integer completion}
      \If{$\ell = \textsc{none}$} \Comment{everything well-defined: a full solution}
        \State \Return \textsc{Feasible}(current assignment)
      \EndIf
      \State \Call{Assert}{$\ell$} \Comment{descend one decision level}
    \EndIf
  \EndLoop
\EndProcedure
\end{algorithmic}
\end{algorithm}

% Source (VERIFIED): ortools/sat/optimization.cc:259
%   RestrictObjectiveDomainWithBinarySearch -- assumption =
%   GetOrCreateAssociatedLiteral(LowerOrEqual(objective_var, target)),
%   ResetAndSolveIntegerProblem({assumption}); on FEASIBLE restrict ub to the
%   solution's objective, on ASSUMPTIONS_UNSAT Enqueue(GreaterOrEqual(obj,target+1)).
\begin{algorithm}[htbp]
\caption{Optimization by bisection: \textsc{MinimizeByBisection} drives \textsc{Solve}}
\label{alg:bisection}
\begin{algorithmic}[1]
\Procedure{MinimizeByBisection}{$\code{obj}$}
  \State $\text{lo} \gets \text{lb}(\code{obj})$; \ $\text{hi} \gets \text{ub}(\code{obj})$; \ \textit{incumbent} $\gets \textsc{none}$
  \While{$\text{lo} < \text{hi}$}
    \State $m \gets \lfloor (\text{lo} + \text{hi})/2 \rfloor$
    \State $r \gets \Call{Solve}{\{\bl{\code{obj} \le m}\}}$ \Comment{assume the trial bound for one call}
    \If{$r = \textsc{Infeasible}$}
      \State $\text{lo} \gets m + 1$ \Comment{the returned core proves $\code{obj} > m$}
    \Else
      \State \textit{incumbent} $\gets r$; \ $\text{hi} \gets \code{obj}(\textit{incumbent})$ \Comment{a solution of cost $\le m$}
    \EndIf
  \EndWhile
  \State \Return \textit{incumbent} as \textsc{Optimal}
\EndProcedure
\end{algorithmic}
\end{algorithm}

The two procedures factor the engine cleanly. \textsc{Solve}'s contract is
feasibility: it returns a feasible assignment or a proof none exists, and it never
manages the objective bound --- that bookkeeping is entirely the wrapper's. The
loop itself carries no optimization logic; the only cost-awareness inside it lives
in the decision strategies, which can be set up, explicitly or implicitly, to
branch toward cheaper solutions (\cref{sec:search}). \textsc{MinimizeByBisection}
(\srcla{optimization.cc}{259}{RestrictObjectiveDomainWithBinarySearch}) keeps a
lower and an upper bound on the
optimal cost and, each round, asks \textsc{Solve} whether a solution of cost at most
the midpoint $m$ exists --- forcing the order-encoding literal
$\bl{\, \code{obj} \le m \,}$ (minted on demand by
\srcla{integer.h}{193}{GetOrCreateAssociatedLiteral}) as an \emph{assumption}, a fact held
only for that one call. This works because the objective is just one more bounded
integer variable (\cref{sec:opt}) and a bound \emph{is} a literal
(\cref{part:foundations}). The two answers move opposite ends of the gap: a solution
pulls the upper bound down to its cost, while an \textsc{Infeasible} hands back a
\emph{core} --- the subset of the assumptions that cannot hold together --- proving
$\code{obj} > m$ and lifting the lower bound. When the two bounds meet, the incumbent
is optimal, and that final infeasibility proof \emph{is} the proof of optimality.
This is the first place \textsc{Solve}'s assumption input and core output earn their
keep; \cref{sec:opt} puts the same $A$ parameter to work in the other
strategies that drive this loop.

This single loop hides a deeper design decision: how the CDCL clause-learning
machinery is coupled to finite-domain propagation. There are broadly two ways to
wire the two together. One option is to keep a complete CDCL \textbf{SAT solver in
charge} and let it own the whole search stack: the trail, the decisions, unit
propagation, conflict analysis, and backjumping. The finite-domain side then talks
to it only through the clause interface, injecting fresh clauses (and the Boolean
variables they need) whenever it notices the solver has missed something a
constraint implies. The SAT engine runs the show; the constraint reasoning feeds
it from outside.

CP-SAT takes the other option: it \textbf{drives the search itself}. It runs its
own propagation, conflict analysis, and search as one integrated core, and the
CDCL clause-learning machinery is a mechanism inside that core rather than an
external engine being fed from the side. This choice runs through every line of
the loop. Boolean literals and integer bounds share one decision-level timeline across two
coupled trails (\cref{sec:trail}), so the \textsc{PropagateToFixedPoint} step runs the CP
propagators directly and a finite-domain propagator pushes an integer-bound fact
onto the integer trail itself (\cref{sec:propstep}); conflict analysis then resolves
over those integer literals natively, and the Boolean clauses and variables are
minted lazily, only at the thin interface where the two worlds meet. This is what
lets the costly apparatus of \cref{part:learning} pay off: a decision and the
reason later learned from it are written in the very same literals, so every
failure feeds straight back into the search that produced it, with no re-encoding
round trip.

The interesting work inside \textsc{Solve}'s conflict case is split between
\textbf{learning} and \textbf{backjumping}. Branching is not a single rule but a walk down an ordered
\textbf{decision strategy} $\mathcal{S}$, a list of heuristics: \textsc{NextDecision}
tries each in turn and returns the first literal one produces, which the loop then
asserts. By default $\mathcal{S}$
holds two such heuristics (\cref{sec:search}): the SAT decision heuristic
(\textbf{VSIDS}), which keeps branching on existing bound literals until every one
has a value, and \textbf{integer completion}, a generic fallback that takes the
first variable still spanning a range and fixes it toward its minimum. Running only
after VSIDS has settled every Boolean, completion in practice sees just the integers
the skeleton has left loose, minting one bound at a time until each domain
collapses to a single value. The loop reports a solution only when \emph{every} heuristic in
$\mathcal{S}$ comes up empty, which is why the leaf test is stricter than ``all
Booleans assigned'': it holds only once every domain has also shrunk to a
singleton.

Which heuristics $\mathcal{S}$ contains, and in what order, is part of the
configuration rather than a fixed feature of the loop. \cref{sec:search}
takes the strategy apart: what each of the two default heuristics does, and how
the parallel portfolio reorders and extends them across workers.

% Source (VERIFIED): the "Search stats" table is built in ortools/sat/stat_tables.cc:84-92
%   (header "...Branches..."; FormatCounter(r.num_branches())). The response field
%   "branches:" (cp_model_solver.cc:713) is NOT a portfolio sum: it is one worker's
%   num_branches(), copied by the statistics postprocessor (see sec:portfolio).
% Empirical (checked against real CP-SAT logs): per-worker Branches has NO fixed
%   relation to the variable count -- orders of magnitude ABOVE it (hard small models,
%   restart/backjump re-descents) or far BELOW it (large models solved at the root).
%   In xl_20x6_120parties.log the response showed branches:1072 (below 2730 vars)
%   because the solution-holding worker barely searched.
\begin{platypuslog}[htb]
Each pass through the \textsc{Assert} line is counted as a \code{Branch} in the
per-worker \textsf{Search stats} table printed at the end of a run. The count has
no fixed relation to the number of variables: it tallies decisions over the
\emph{whole} search, so restarts and backjumping can drive a searching worker to
many times the variable count, while a worker that rides presolve and root
propagation can finish with far fewer. Within one worker's row it sits well below
the two propagation columns beside it (\cref{sec:propstep}): every decision sets
off a full cascade, so \code{BoolPropag} and \code{IntegerPropag} both run far
ahead of \code{Branches} wherever real search happens. Read these per worker; the
response summary's \code{branches:} field reports only the worker that returned
the solution (\cref{sec:portfolio}), which can be far smaller.
\end{platypuslog}

\textsc{Solve} and its optimizer, together with strong propagators, are the search
core. Bisection is one clean way to drive the feasibility loop;
\cref{sec:opt} surveys the others, including the even simpler linear scan and
the core-guided strategies the portfolio runs.

\subsection{Inside ``propagate to a fixed point''}\label{sec:propstep}
% Sources (VERIFIED):
%   ortools/sat/sat_solver.cc -- SatSolver::Propagate(): loops over an ordered
%     propagators_ list; the inner loop BREAKS as soon as the trail grows, so it
%     restarts from the cheapest propagator. Registration order: binary_implication_graph_,
%     clauses_propagator_, enforcement_propagator_, pb_constraints_, external,
%     then GenericLiteralWatcher (last_propagator_). The in-source comment states
%     the intent ("loop over and test again the highest priority constraint types").
%   ortools/sat/clause.cc -- BinaryImplicationGraph::Propagate() and
%     ClauseManager::Propagate() ARE the unit-propagation engines (the "first row").
%   ortools/sat/integer.cc -- GenericLiteralWatcher::Propagate(): priority buckets
%     queue_by_priority_ (default priority 1; 0 = cheap), drains low priority first,
%     resets to priority 0 when an integer bound is enqueued, RETURNS to the SAT
%     loop when a Boolean literal is enqueued.
%   ortools/sat/integer.h -- WatchLiteral/WatchLowerBound/WatchUpperBound +
%     PropagatorInterface::IncrementalPropagate(watch_indices): change-driven waking.

The loop of \cref{sec:loop} hides its most important line:
``propagate to a fixed point.'' What actually happens there, and where does the
SAT engine sit relative to the CP propagators?

\begin{platypuswarning}[htb][Not generate-and-check]
CP-SAT does \emph{not} let the SAT solver guess a complete Boolean assignment
and then hand it to the CP propagators for checking, as in generate-and-check
lazy SMT or CEGAR. Propagation is \textbf{online}: it runs after \emph{every}
decision, before the next one, and the propagators always operate on the current
\emph{partial} state, namely the bounds as they stand, never on a finished
assignment.
\end{platypuswarning}

\textbf{A layered, cheapest-first cascade.} Propagators are kept in a fixed
order, with the cheapest and most prolific first, and the engine runs them as a
chain:

\begin{enumerate}[leftmargin=1.6em,topsep=0.3em,itemsep=0.25em]
  \item the \textbf{binary-implication graph}: implications from two-literal
  clauses, the cheapest form of unit propagation;
  \item the \textbf{clause manager}: unit propagation on general clauses via the
  two-watched-literals scheme;
  \item enforcement and \textbf{pseudo-Boolean} constraints;
  \item the \textbf{\srcla{integer.h}{1324}{GenericLiteralWatcher}}, which fans
  out to all integer and CP propagators: linear,
  \srcla{all\_different.h}{52}{AllDifferent},
  \srcla{disjunctive.h}{49}{NoOverlap},
  \srcla{cumulative.h}{48}{Cumulative}, and so on.
\end{enumerate}

The first two layers \emph{are} classical SAT unit propagation. This is the
``first-row propagator'' picture: cheap, always-run Boolean reasoning fires
before any expensive global constraint is even consulted.

\textbf{Restart-on-progress is the key scheduling rule.} The instant any
propagator deduces something new, by enqueuing a literal or tightening a bound,
the chain breaks and starts again \emph{from the top}. Thus, when an expensive
global constraint such as \srcla{cumulative.h}{48}{Cumulative} tightens a bound,
control immediately returns to binary-implication and clause unit propagation.
Those cheap layers then run to their own fixed point again before any expensive
propagator is reconsidered. Cheap reasoning is always exhausted before the
engine spends time on expensive reasoning.

The same discipline repeats one level down inside
\srcla{integer.h}{1324}{GenericLiteralWatcher}. CP propagators are themselves
bucketed by priority: priority 0 for cheap propagators such as linear and
scheduling helpers, and higher priorities for more expensive propagators such as
\srcla{all\_different.h}{52}{AllDifferent} or
\srcla{cumulative.h}{48}{Cumulative}. The watcher drains priority 0 before
priority 1, but resets to priority 0 as soon as a cheap propagator tightens a
bound. If a propagator pushes a \emph{Boolean} literal, control returns to the
SAT layer immediately, so unit propagation re-runs before anything costlier.

\textbf{Change-driven waking.} None of this rescans every constraint in every
round. At registration, a propagator declares exactly which literals and
variable bounds it cares about (\srcla{integer.h}{1369}{WatchLiteral}, \srcla{integer.h}{1370}{WatchLowerBound},
\srcla{integer.h}{1371}{WatchUpperBound}). It is then woken only when one of \emph{those} facts
changes, and it receives the changed watch indices through
\srcla{integer.h}{1299}{IncrementalPropagate} rather than recomputing from scratch. The fixed-point
loop is therefore sparse: only the constraints touched by the last deduction are
re-run, and they are usually told precisely what moved.

Putting this together, a single decision triggers a cascade that first exhausts
the cheapest propagation, escalates to expensive global constraints only when
the cheap layers stall, returns to cheap propagation the moment anything new is
deduced, and wakes only constraints whose watched facts changed. Only when the
\emph{entire} cascade stalls does control return to branching; if some
propagator instead empties a domain, control goes to the conflict analysis of
\cref{sec:learning}.

% Source (VERIFIED): "Search stats" columns "BoolPropag"/"IntegerPropag" =
%   FormatCounter(r.num_binary_propagations())/FormatCounter(r.num_integer_propagations())
%   in ortools/sat/stat_tables.cc:96-97. The two counters come from two structures:
%   num_binary_propagations() := SatSolver::num_propagations() (Boolean trail, binary
%   + clause unit propagation), num_integer_propagations() := IntegerTrail::num_enqueues()
%   (integer trail). Naming trap: the response field is PRINTED as "propagations:" but
%   is num_binary_propagations() (cp_model_solver.cc:715-718 comment says so), the
%   Boolean count only; the response prints just ONE worker's figures (see sec:portfolio).
% Empirical (real CP-SAT logs): do NOT claim IntegerPropag > BoolPropag -- the
%   order between them flips by model (clause-heavy encodings make BoolPropag the
%   larger, e.g. 65.7M vs 36.9M). What IS robust: both > Branches in every
%   searching worker. So anchor the message on the propagation-vs-Branches gap.
\begin{platypuslog}[htb]
The two trails are tallied by two separate columns of the \textsf{Search stats}
table. \code{BoolPropag} counts propagations on the Boolean trail (the
binary-implication and clause layers), read from the SAT solver's own counter;
\code{IntegerPropag} counts bound tightenings pushed onto the integer trail, read
from that trail's enqueue count. In a worker that genuinely searches they are
usually the two largest columns by a wide margin, the quantitative face of
``propagation does most of the work, branching only a little.'' Which of the two
leads is model-dependent: clause-heavy encodings can put \code{BoolPropag} ahead,
so the message lies in their size relative to \code{Branches}, not in the split
between them. The response summary echoes the pair as \code{integer\_propagations:}
and --- a naming trap --- \code{propagations:} for the Boolean count, but, like the
other summary counters, only for the single worker that returned the solution
(\cref{sec:portfolio}).
\end{platypuslog}

\subsection{Inside ``search'': choosing the next decision}\label{sec:search}
% Sources (VERIFIED):
%   ortools/sat/sat_decision.h:58 -- VSIDS-style decaying activity, phase saving.
%   ortools/sat/integer_search.cc:422 SatSolverHeuristic -- all_assigned test
%     (returns no decision when every Boolean is assigned);
%   integer_search.cc:61 AtMinValue (min-first completion);
%   integer_search.cc:80 GreaterOrEqualToMiddleValue (bisection);
%   integer_search.cc:119 SplitAroundLpValue (LP value);
%   integer_search.cc:1224 SequentialSearch({SatSolverHeuristic, fixed_search})
%     -- completion (fixed_search) is the last heuristic in the default policy.

The companion to ``propagate to a fixed point'' is the other line the loop hides:
how the next decision is actually chosen. The integrated design of
\cref{sec:loop} --- one decision-level timeline across its two trails, with clause learning a mechanism
inside the core --- settles how decisions are represented and learned from, but
not \emph{which} heuristic picks the next one. That is the configurable decision
strategy $\mathcal{S}$ of \cref{sec:loop}: the default leads with SAT
branching and falls back to finite-domain branching, but other workers reorder
these or lead with an LP or a fixed integer order instead. The two heuristics
below are the entries of that default $\mathcal{S}$, and the building blocks the
other strategies recombine.

\textbf{The SAT decision heuristic.} This entry branches on a bound literal that
already exists. CP-SAT's SAT core chooses it by \textbf{VSIDS}
(\srcl{sat\_decision.h}{58}): it keeps a decaying activity score per variable,
bumps the variables that appear in each conflict, and branches on whichever has
been most involved in recent failures, with a saved polarity (\emph{phase
saving}) picking true or false. Because an integer bound is a reified literal
(\cref{sec:bounds}), this one heuristic ranges over Boolean and integer
variables alike: the literals it scores are exactly the bound literals conflict
analysis has been resolving over. This is the sense in which this entry branches
from the SAT side.

\textbf{Integer completion.} VSIDS can branch only on literals that already
exist, and the sparse skeleton (\cref{sec:completion}) lets those run out
while an integer is still ranged: every existing bound for $x$ is assigned,
yet $x \in [3,7]$ because no bound between $3$ and $7$ was ever minted. When
the SAT heuristic reports nothing left to decide, every Boolean is assigned
(\srcla{integer\_search.cc}{422}{SatSolverHeuristic}), and only then does an
\emph{integer-completion} heuristic step in: it picks a still-ranged variable,
mints a fresh bound for it, and branches on that. The default composition puts
completion last for exactly this reason (\srcl{integer\_search.cc}{1224}):
exhaust the existing Booleans first, and fall back to minting a new one only when
the skeleton is settled but some integer is still loose.

Which bound to mint is a value-selection choice. The default branches the
variable toward its current minimum, asserting $\bl{x \le \text{lb}(x)}$ (\srcla{integer\_search.cc}{61}{AtMinValue}): the first guess tries
the smallest surviving value, and the other branch pushes the lower bound up.
Other configurations bisect instead, minting a middle bound
$\bl{x \ge \text{lb} + \lceil(\text{ub}-\text{lb})/2\rceil}$
(\srcla{integer\_search.cc}{80}{GreaterOrEqualToMiddleValue}) to halve the range
per decision, or split around the current LP value. Each is just a rule for which
not-yet-minted bound becomes the next decision.

This is why search terminates on $\text{lb} = \text{ub}$ rather than when the
Booleans are exhausted: completion keeps minting bounds for an unfixed
variable instead of declaring victory once the SAT skeleton is satisfied. The
propagation fixed point of \cref{sec:propstep} becomes a \emph{solution}
only when every integer has been pinned to a single value, one missing literal at
a time.

\textbf{How workers vary the strategy.} These two heuristics are only the default
$\mathcal{S}$. Nothing in the loop of \cref{sec:loop} fixes \emph{which}
strategy it runs, and the parallel portfolio (\cref{sec:portfolio})
exploits this, running workers whose $\mathcal{S}$ differs in both contents and
order, selected through the \code{search\_branching} parameter:

\begin{itemize}[leftmargin=1.6em,topsep=0.3em,itemsep=0.25em]
  \item \code{auto} (\srcla{integer\_search.cc}{1224}{AUTOMATIC\_SEARCH}, the default):
  $\langle\text{VSIDS},\ \text{integer completion}\rangle$;
  \item \code{fixed} (\srcla{integer\_search.cc}{1239}{FIXED\_SEARCH}):
  $\langle\text{fixed integer order},\ \text{VSIDS}\rangle$, the integer-style
  decisions first, with VSIDS only settling any leftover Booleans;
  \item \code{reduced\_costs} (\srcla{integer\_search.cc}{1288}{LP\_SEARCH}):
  $\langle\text{LP reduced-cost},\ \text{VSIDS},\ \text{integer completion}\rangle$;
  \item \code{pseudo\_costs} (\srcla{integer\_search.cc}{1310}{PSEUDO\_COST\_SEARCH}):
  $\langle\text{pseudo-cost},\ \text{VSIDS},\ \text{integer completion}\rangle$;
  \item \code{portfolio} and \code{quick\_restart}
  (\srcla{integer\_search.cc}{1282}{PORTFOLIO\_SEARCH}, \srcla{integer\_search.cc}{1319}{PORTFOLIO\_WITH\_QUICK\_RESTART\_SEARCH}):
  a randomized rotation over such lists, reshuffled on restart.
\end{itemize}

Only the contents and order of $\mathcal{S}$ change between workers; the loop
around it does not.

% Sources (VERIFIED against sat_parameters.proto / cp_model.proto):
%   search_branching (SearchBranching enum) field 82, default AUTOMATIC_SEARCH:
%     sat_parameters.proto:1194 (field), enum 1151-1193. FIXED_SEARCH=1 (proto:1161).
%   The user-supplied fixed order is the repeated search_strategy field of the
%     CpModelProto: cp_model.proto:654 (repeated DecisionStrategyProto), message
%     DecisionStrategyProto cp_model.proto:520, VariableSelectionStrategy enum
%     :539 (CHOOSE_FIRST=0 :540, CHOOSE_LOWEST_MIN=1 :541, ...),
%     DomainReductionStrategy enum :552 (SELECT_MIN_VALUE=0 :553,
%     SELECT_MAX_VALUE=1 :554, ...). Python API: model.add_decision_strategy().
%   Caveat re portfolio: a single global search_branching binds only the workers
%     that honour it; see the subsolver-restriction box in \cref{sec:portfolio}.
\begin{platypusparam}[htb]
You can also supply your own search strategy $\mathcal{S}$ instead of leaving it to the
portfolio. The \srcla{sat\_parameters.proto}{1194}{search\_branching} parameter (default \code{AUTOMATIC\_SEARCH}) selects the policy,
and \code{FIXED\_SEARCH} makes the solver follow an order \emph{you} specify,
walked by \srcla{cp\_model\_search.cc}{210}{ConstructUserSearchStrategy}.
That order is a per-model object: each
\srcla{cp\_model.proto}{520}{DecisionStrategyProto} (the repeated
\srcla{cp\_model.proto}{654}{search\_strategy} field, exposed as
\code{add\_decision\_strategy} in the Python API) pairs a
\textbf{variable-selection} rule (\code{CHOOSE\_FIRST},
\code{CHOOSE\_LOWEST\_MIN}, \dots) with a \textbf{value} rule
(\code{SELECT\_MIN\_VALUE}, \code{SELECT\_MAX\_VALUE}, \dots), letting you
hand-encode a custom variable and value order rather than leaving the choice to
the automatic activity heuristic.
\end{platypusparam}

\subsection{A small worked example}

The following tiny model illustrates the control loop without introducing an
objective yet:
\[
\begin{aligned}
  &a, b, c \in \{0, 1\} \quad\text{(Booleans)}, \qquad x \in [0, 5] \\
  &\text{C1:}\ a + b + c \ge 2 \quad\text{(at least two of } a,b,c \text{ are true)}\\
  &\text{C2:}\ x \ge 3a \\
  &\text{C3:}\ x \le 2
\end{aligned}
\]

\textbf{Root propagation.} C3 tightens $x$ to $x \le 2$. C2 says
$x \ge 3a$; in combination with $x \le 2$, it forces $a = 0$, because $a = 1$
would require $x \ge 3$. Now C1 becomes $b + c \ge 2$, forcing $b = 1$ and
$c = 1$. All of this is pure deduction, with no decisions at all.

\textbf{Completion search.} At this point the state is consistent, but not yet a
solution. The Boolean variables are fixed, whereas $x$ still ranges over
$[0,2]$: once $a=0$, no remaining constraint distinguishes $x=0$, $x=1$, and
$x=2$. Recall \cref{sec:search}: search terminates only when every
variable has a singleton domain. The completion brancher therefore still takes a
decision on $x$. Under the minimum-first completion heuristic, it tries
$x \le 0$; together with the existing lower bound $x \ge 0$, this fixes $x=0$.
Only then is the total assignment $(a,b,c,x) = (0,1,1,0)$ reported.

The example is intentionally conflict-free. Its purpose is to show the ordinary
case in which propagation and branching cooperate: propagation performs the
forced deductions, and search supplies the remaining choices needed to turn a
consistent state into a complete assignment. If a later decision in another
branch did lead to a conflict, the same reasons attached to these propagations
would feed the conflict analysis described in \cref{sec:learning}.

\subsection{Optimization: solving feasibility, repeatedly}\label{sec:opt}
% Sources (VERIFIED): ortools/sat/optimization.cc
%   Linear-scan path: on FEASIBLE, objective = integer_trail->LowerBound(objective_var),
%   then Enqueue(IntegerLiteral::LowerOrEqual(objective_var, objective - 1)) and
%   re-solve -> next solution must be strictly better. (Other strategies exist,
%   e.g. core-based CoreBasedOptimizer; the bound-tightening idea is the shared core.)

Optimization treats the objective as one more bounded integer variable. This is
literal, not a metaphor: the model holds a dedicated integer variable,
\srcla{cp\_model\_mapping.h}{41}{objective\_var}, constrained to equal the
objective expression, and to
minimize is to drive down its upper bound. Because it is an ordinary integer
variable, it sits on the integer trail, propagators watch and tighten it, and
parallel workers share improved bounds on it like any other bound
(\cref{sec:portfolio}). This is what lets \textsc{MinimizeByBisection}
(\cref{alg:bisection}) optimize with no optimization-specific search: it
only ever moves a bound on \code{objective\_var} and re-runs the same feasibility
loop. Bisection is one way to move that bound; it is not the only one.

\textbf{Linear scan} (\srcl{optimization.cc}{212}) is even simpler, and it is the
one strategy that needs no assumptions at all. Find a solution of cost $C$, then forbid that cost and
everything worse and search again. Because the objective only ever has to improve,
the cutoff is \emph{monotone}: it never has to be taken back. So rather than
assuming it for a single call, \textsc{MinimizeByLinearScan} \emph{permanently}
tightens \code{objective\_var}'s upper bound to $C - 1$ at the root
(\cref{alg:linscan}); the bound becomes part of the model, propagates like
any other, and every later solution is strictly cheaper. When the tightened model
is infeasible, the last $C$ was optimal --- again the infeasibility proof \emph{is}
the proof of optimality. This is the contrast with bisection that matters: bisection
probes a bound it may have to retract, so it \emph{assumes} it; linear scan tightens
a bound it will never relax, so it \emph{asserts} it for good. Either way, CP-SAT
reports both an \textbf{incumbent} (the best feasible cost so far) and a
\textbf{bound} (the best cost ruled out below); when the two meet the gap is zero
and optimality is proved, and good incumbents, propagating like any other bound,
accelerate that proof.

% Source (VERIFIED): ortools/sat/integer_search.cc -- objective-directed branching,
%   e.g. SplitAroundLpValue (119) descends toward the LP relaxation's values; the
%   solution Solve returns is therefore typically well below the posted cutoff, so
%   the realized step exceeds the single unit the cutoff demands.
\begin{platypusinfo}[htb][The step size comes from the search, not the cutoff]
The cutoff posted after each solution only demands that the \emph{next} one be
strictly better: it tightens $\bl{\, \code{obj} \le \code{obj}(\textit{incumbent}) - 1 \,}$,
a single unit. How far each step actually moves is set by the search, not by that
cutoff. Objective-directed decision strategies, such as branching toward the LP
relaxation's values (\cref{sec:search}), steer each \textsc{Solve} call toward a
good assignment rather than merely a feasible one just under the cutoff, so the
solution returned is usually \emph{much} cheaper than one unit better. Linear scan
therefore tends to descend in large, uneven jumps rather than crawling unit by
unit. The same effect helps bisection: it banks the returned incumbent's true cost
(\cref{alg:bisection}, $\text{hi} \gets \code{obj}(\textit{incumbent})$), which
often sits well below the probed midpoint.
\end{platypusinfo}

% Source (VERIFIED): ortools/sat/optimization.cc -- linear-scan path: on FEASIBLE,
%   Enqueue(LowerOrEqual(objective_var, objective - 1)) at level 0 and re-solve.
\begin{algorithm}[htbp]
\caption{Linear scan: \textsc{MinimizeByLinearScan} tightens a monotone bound}
\label{alg:linscan}
\begin{algorithmic}[1]
\Procedure{MinimizeByLinearScan}{$\code{obj}$}
  \State \textit{incumbent} $\gets \textsc{none}$
  \Loop
    \State $r \gets \Call{Solve}{\emptyset}$ \Comment{search under the cutoffs posted so far}
    \If{$r = \textsc{Infeasible}$}
      \State \Return \textit{incumbent} as \textsc{Optimal} if one was found, else \textsc{Infeasible}
    \EndIf
    \State \textit{incumbent} $\gets r$ \Comment{a strictly better solution; record it}
    \LineComment{No assumption needed: the cutoff only ever tightens, so post it permanently at the root}
    \State permanently tighten $\bl{\, \code{obj} \le \code{obj}(\textit{incumbent}) - 1 \,}$ at the root
  \EndLoop
\EndProcedure
\end{algorithmic}
\end{algorithm}

\textbf{Core-guided} workers
(\srcla{optimization.cc}{444}{CoreBasedOptimizer}) push the same machinery
further. Instead of bounding
\code{obj} as a whole, they assume the individual objective terms at their best
values and call \textsc{Solve}; each returned core
(\srcla{optimization.cc}{376}{FindCores}) is a set of terms that cannot
all stay at their best, which they use to reformulate the objective and lift the
lower bound (sometimes minimizing the core first). Bound-focused workers
(\code{objective\_lb\_search}, \code{objective\_shaving}) specialize in the dual
side. They differ in \emph{how} they move the incumbent and the bound, but every
one of them is an outer loop over the single \textsc{Solve} of
\cref{alg:loop}: they share its propagation, learning, and search, and
vary only what they assert or assume on the objective between calls.

% Source (VERIFIED): ProgressMessage()/SatProgressMessage() in
%   ortools/sat/cp_model_solver_logging.cc:43-92 -- "#%-5s %6.2fs best:... next:[lb,ub]";
%   solution counter "#1,#2,..." (num_solutions_), throttled "#Bound", and "#Done".
\begin{platypuslog}[htb]
The optimizer's bound-tightening is what generates the streaming progress lines: each
improving incumbent prints as \code{\#1}, \code{\#2}, \ldots, with its cost as
\code{best:} and the still-open gap as \code{next:[lb,ub]}. These lines are
interleaved with throttled \code{\#Bound} lines whenever a worker improves the
proof bound. When \code{best} meets the bound, the gap is zero and a final
\code{\#Done} line closes the search.
\end{platypuslog}
